% =============================================================================
%  2-Desarrollo.tex — Cláusulas del acuerdo
% =============================================================================

\section{Cláusulas}

% ── Primera ───────────────────────────────────────────────────────────────────
\clausula{Primera}{Objeto}

El presente acuerdo tiene por objeto establecer las reglas mediante las cuales cualquiera de Las Partes puede decir que no o detener una conversación o discusión que la desborde, así como la forma y los plazos para retomar el tema de manera ordenada y respetuosa.

% ── Segunda ───────────────────────────────────────────────────────────────────
\clausula{Segunda}{Derecho a decir no}

\begin{enumerate}[label=\arabic*., leftmargin=*, itemsep=0.3em]
  \item Cualquiera de Las Partes puede decir que no, expresando que eso es lo que necesita, sin estar obligada a explicar sus sentimientos ni a justificar su negativa.
  \item La otra Parte debe respetar la negativa, sin insistir, presionar ni exigir justificación en ese momento.
\end{enumerate}

% ── Tercera ───────────────────────────────────────────────────────────────────
\clausula{Tercera}{Opinión clara y motivo de la conversación}

\begin{enumerate}[label=\arabic*., leftmargin=*, itemsep=0.3em]
  \item Quien inicie o plantee un tema debe expresar con claridad su opinión y la razón por la cual desea hablarlo.
  \item Esta obligación no contradice la Cláusula Segunda: el deber de explicar recae en quien abre el tema, no en quien decide decir que no o detener la conversación.
\end{enumerate}

% ── Cuarta ────────────────────────────────────────────────────────────────────
\clausula{Cuarta}{Pausa activa y frases clave}

\begin{enumerate}[label=\arabic*., leftmargin=*, itemsep=0.3em]
  \item \textbf{Frases clave.} La pausa activa se solicita al pronunciar cualquiera de las siguientes frases:
  \begin{itemize}[itemsep=0.1em]
    \item \fraseclave{No quiero seguir con la conversación.}
    \item \fraseclave{No quiero hablar sobre ello.}
  \end{itemize}

  \item \textbf{Efecto inmediato.} Al pronunciarse una frase clave, la conversación se detiene en ese momento. La Parte receptora no debe continuar con el tema, agregar un último argumento ni responder con reproches.

  \item \textbf{Modalidades.} Al solicitar la pausa, la Parte solicitante puede elegir una de dos modalidades:
  \begin{enumerate}[label=\alph*), itemsep=0.2em]
    \item \textbf{Por tiempo:} especificando un lapso que no podrá ser mayor a \textbf{8 horas}.
    \item \textbf{Por circunstancia:} la pausa dura hasta que la Parte solicitante recupere su estabilidad, sin un plazo fijo.
  \end{enumerate}

  \item \textbf{Fin de la pausa por circunstancia.} La pausa por circunstancia termina cuando la estabilidad de la Parte solicitante es visible para ambas Partes frente a los estímulos físicos, sensoriales o emocionales que la desbordaron. Se considera estabilidad visible, entre otras, la regulación:
  \begin{enumerate}[label=\alph*), itemsep=0.2em]
    \item \textbf{Física:} por ejemplo, de la temperatura corporal.
    \item \textbf{Emocional:} por ejemplo, ya no hay llanto, gritos ni desborde.
    \item \textbf{Social:} disposición a convivir y a comunicarse con la otra Parte.
  \end{enumerate}

  \item \textbf{Naturaleza de la pausa.} La pausa activa es una suspensión de la conversación, no de la relación. No implica distanciamiento ni huida sentimental, ni puede utilizarse como castigo o como forma de ignorar a la otra Parte.

  \item \textbf{Obligación de retomar la interacción.} Concluida la pausa, la Parte solicitante está obligada a retomar la interacción con la otra Parte. Retomar la interacción no significa retomar la conversación del tema, la cual se rige por la Cláusula Sexta. Al retomar la interacción, la Parte solicitante declara que es consciente de que se encuentra regulada emocional y sensorialmente para mantener una interacción sana.

  \item \textbf{Trato respetuoso.} Durante la pausa y al retomar la interacción, Las Partes deben tratarse con respeto: sin lenguaje ofensivo, indirectas ni expresiones similares.

  \item \textbf{Reiteración de la pausa.} La pausa activa puede solicitarse las veces que sea necesario, sin límite, incluso durante una conversación de resolución.

  \item \textbf{Pausa de la Parte receptora.} Si al retomar la interacción la Parte solicitante no se encuentra regulada y la interacción no es sana, la Parte que estaba en espera puede solicitar a su vez una pausa activa para proteger la integridad de ambas Partes y la sana interacción. Esta pausa no se considera represalia y se rige por las mismas reglas de esta cláusula.
\end{enumerate}

% ── Quinta ────────────────────────────────────────────────────────────────────
\clausula{Quinta}{Excepciones por integridad, seguridad o salud}

\begin{enumerate}[label=\arabic*., leftmargin=*, itemsep=0.3em]
  \item La pausa activa queda suspendida cuando la integridad, la seguridad o la salud de cualquiera de Las Partes se encuentre en riesgo. De manera enunciativa, mas no limitativa, se consideran los siguientes casos:
  \begin{enumerate}[label=\alph*), itemsep=0.2em]
    \item Dolor físico.
    \item Dolor emocional extremo, manifestado en desmayos, llanto o gritos.
    \item Agresión verbal.
    \item Agresión física, desde el contacto no deseado hasta el contacto agresivo o cualquier contacto que vulnere la integridad física.
  \end{enumerate}

  \item En estos casos, la prioridad es la seguridad y el bienestar de Las Partes: atender a quien lo necesite y, de ser necesario, pedir ayuda a terceros o a servicios de emergencia. Una frase clave no puede usarse para impedir que una de Las Partes brinde o pida ayuda.

  \item Ante una agresión verbal o física, cualquiera de Las Partes puede separarse físicamente de la otra. Esta separación no se considera distanciamiento ni incumplimiento de la Cláusula Cuarta.
\end{enumerate}

% ── Sexta ─────────────────────────────────────────────────────────────────────
\clausula{Sexta}{Plazos para retomar el tema}

\begin{enumerate}[label=\arabic*., leftmargin=*, itemsep=0.3em]
  \item \textbf{Plazo mínimo.} El tema que originó la pausa activa solo puede retomarse cuando hayan transcurrido al menos \textbf{24 horas} desde el suceso.

  \item \textbf{Forma.} El tema solo puede retomarse si una de Las Partes inicia una conversación de resolución, es decir, orientada a resolver, comprender o entender la causa de la pausa activa, y no a reprochar ni a reabrir la discusión en los mismos términos.

  \item \textbf{Plazo máximo.} El tema puede retomarse a más tardar \textbf{una semana} (7 días naturales) después del suceso.

  \item \textbf{Conflicto solucionado.} Si ninguna de Las Partes retoma el tema dentro del plazo máximo, este se considerará \textbf{\enquote{Conflicto solucionado}}.

  \item \textbf{Efecto.} Un conflicto solucionado no podrá ser usado por ninguna de Las Partes como argumento, reproche o antecedente en conflictos futuros.

  \item \textbf{Nueva pausa sobre el mismo tema.} Si durante una conversación de resolución, o al retomar la interacción, se solicita una nueva pausa activa sobre el mismo tema, esta constituye un nuevo suceso y los plazos mínimo y máximo de esta cláusula vuelven a contarse desde ese momento.
\end{enumerate}

\begin{table}[H]
  \centering
  \small
  \begin{tabularx}{\textwidth}{@{} l >{\raggedright\arraybackslash}p{3.4cm} X @{}}
    \toprule
    \textbf{Momento}         & \textbf{Plazo}        & \textbf{Regla} \\
    \midrule
    Pausa por tiempo         & Máximo 8 horas        & La Parte solicitante retoma la interacción, no el tema. \\
    Pausa por circunstancia  & Hasta que haya estabilidad visible & La Parte solicitante retoma la interacción, no el tema. \\
    Retomar el tema          & Mínimo 24 horas       & Solo mediante una conversación de resolución. \\
    Cierre del tema          & Máximo 7 días         & Si no se retoma: \enquote{Conflicto solucionado}. \\
    Nueva pausa              & Sin límite            & Es un nuevo suceso: los plazos de 24 horas y 7 días se reinician. \\
    \bottomrule
  \end{tabularx}
  \caption{Resumen de plazos, contados a partir del suceso.}
  \label{tab:plazos}
\end{table}

% ── Séptima ───────────────────────────────────────────────────────────────────
\clausula{Séptima}{Vigencia y modificaciones}

\begin{enumerate}[label=\arabic*., leftmargin=*, itemsep=0.3em]
  \item El presente acuerdo entra en vigor en la fecha de su firma y tiene vigencia indefinida.
  \item Cualquiera de Las Partes puede solicitar su revisión. Toda modificación debe hacerse de mutuo acuerdo y por escrito, mediante una adenda o un nuevo acuerdo.
\end{enumerate}
